\naslov{Izidi, dogodki, verjetnosti}

\begin{vo}{Kaj je množica $\Omega$ vseh možnih izidov? Povej nekaj primerov.}
  To je množica, ki hrani vse možne rezultate nekega poskusa.
  Pri mešanju kupa $n$ kart velja $\Omega = S_n$, pri $n$-kratnem metu kovanca
  je to $\Omega = \{G, S\}^n$, itd.
\end{vo}

\begin{definicija}
  Družina $\mathcal{F}$ podmnožic množice $\Omega$ je \pojem{$\sigma$-algebra},
  če velja:
  \begin{itemize}
  \item $\Omega \in \mathcal{F}$,
  \item $A \in \mathcal{F} \implies A\complement \in \mathcal{F}$,
  \item $A_1, A_2, \ldots \in \mathcal{F} \implies \bigcup_{i} A_i \in \mathcal{F}$.
  \end{itemize}
\end{definicija}

\begin{definicija}
  Naj bo $\Omega$ množica možnih izidov, in $\mathcal{F}$ $\sigma$-algebra nad
  $\Omega$.
  \pojem{Verjetnost} je preslikava $P: \mathcal{F} \to [0,1]$, za katero velja
  $P(\Omega) = 1$, in kjer za disjunktne dogodke $A_1, A_2, \ldots \in
  \mathcal{F}$ velja $P(\bigcup_i A_i) = \sum_i P(A_i)$.
\end{definicija}

\begin{opomba}
  To sta \emph{aksioma Kolmogorova}.
\end{opomba}

\vprasanje{Kaj je verjetnost?}

\begin{izrek}[Formula za vključitve in izključitve]
  Naj bodo $A_1, \ldots, A_n$ dogodki.
  Potem velja
  \[
    P(\bigcup_{i=1}^n A_i) = \sum_{i=1}^n P(A_i) - \sum_{1 \le i < j \le n}
    P(A_i \presek A_j) + \ldots + (-1)^{n-1} P(\bigcap_{i=1}^n A_i).
  \]
\end{izrek}

\begin{proof}
  Definirajmo dogodke
  \[
    B_r = \{\omega \in \Omega \such \text{$\omega$ je vsebovan v natanko $r$
      množicah $A_i$} \}.
  \]
  To so disjunktni dogodki, za katere velja $\bigcup_i A_i = \bigcup_r B_r$.
  Sledi
  \[
    P(\bigcup_{i=1}^n A_i) = \sum_{r=1}^n P(B_r).
  \]
  Poglejmo si, kolikokrat smo v formuli v izreku šteli vsako izmed množic $B_r$.
  Ta množica je vsebovana v preseku do $r$ dogodkov, torej se v prvem členu
  pojavi $r$-krat, v drugem $\binom{r}{2}$, v tretjem $\binom{r}{3}$, itd.
  Vsota je tedaj
  \[
    \binom{r}{1} - \binom{r}{2} + \binom{r}{3} - \ldots + (-1)^r \binom{r}{r} = 1,
  \]
  kar lahko izpeljemo iz razvoja izraza $0 = (1 - 1)^r$.
\end{proof}

\vprasanje{Povej formulo za izključitve in izključitve. Kaj je ideja dokaza?}

\begin{lema}
  Naj bodo $A_1, A_2, \ldots$ dogodki.
  Če je $A_1 \subseteq A_2 \subseteq A_3 \subseteq \cdots$, je verjetnost unije
  \[
    P(\bigcup_{i=1}^\infty A_i) = \lim_{n \to \infty} P(A_n).
  \]
  Če namesto tega velja $A_1 \supseteq A_2 \supseteq A_3 \supseteq \cdots$, je
  \[
    P(\bigcap_{i=1}^\infty A_i) = \lim_{n \to \infty} P(A_n).
  \]
\end{lema}

\begin{proof}
  Druga formula sledi iz De Morganovih pravil, dokažemo samo prvo.
  Zapišemo
  \[
    \bigcup_{i=1}^\infty A_i = A_1 \unija (A_2 \brez A_1) \unija (A_3 \brez (A_1
    \unija A_2)) \unija \ldots
  \]
  To so disjunktni dogodki, torej zanje velja
  \begin{align*}
    P(\bigcup_{i=1}^\infty A_i) &= P(A_1) + \sum_{k=2}^\infty P(A_k \brez (A_1 \unija \ldots \unija A_{k-1})) \\
                                &= \lim_{n \to \infty} (P(A_1) + \sum_{k=2}^n P(A_k \brez (A_1 \unija \ldots \unija A_{k-1}))) \\
                                &= \lim_{n \to \infty} P(\bigcup_{k=1}^n A_k \brez (A_1 \unija \ldots \unija A_{k-1})) \\
                                &= \lim_{n \to \infty} P(A_n).
  \end{align*}
\end{proof}

\begin{lema}[Prva Borel-Cantorjeva lema]
  Naj bodo $A_1, A_2, \ldots$ dogodki, za katere velja $\sum_i P(A_i) < \infty$.
  Definiramo $\cl{A} = \{\omega \in \Omega \such \text{ \(\omega\) je vsebovan v
    neskončno mnogo \(A_k\) }\}$.
  Tedaj velja $P(\cl{A}) = 0$.
\end{lema}

\begin{proof}
  Prepričamo se lahko, da velja $\cl{A} = \bigcap_{n=1}^\infty
  \bigcup_{m=n}^\infty A_m$.
  Te unije so padajoče za $n \to \infty$, zatorej po prejšnji lemi velja
  \[
    P(\cl{A}) = \lim_{n \to \infty} P(\bigcup_{m=n}^\infty A_m).
  \]
  Iz dokaza prešnje leme vidimo, da velja sklep
  \[
    P(\bigcup_{k=1}^n A_k) \le \sum_{k=1}^n P(A_k)
    \implies P(\bigcup_{k=1}^\infty) \le \sum_{k=1}^\infty P(A_k).
  \]
  Torej velja
  \[
    P(\cl{A}) \le \lim_{n \to \infty} \sum_{k=n}^\infty P(A_k).
  \]
  Izraz na desni pa je rep konvergenčne vrste, torej je limita enaka $0$.
\end{proof}

\vprasanje{Povej in dokaži prvo Borel-Cantorjevo lemo.}

\podnaslov{Pogojna verjetnost in neodvisnost}

\begin{definicija}
  Naj bo $B$ dogodek s $P(B) > 0$.
  \pojem{Pogojna verjetnost} dogodka $A$ glede na $B$ je
  \[
    P(A \such B) = \frac{P(A \presek B)}{P(B)}.
  \]
\end{definicija}

\vprasanje{Kaj je pogojna verjetnost?}